% Concise manuscript text; implementation and checks are archived with the analysis.
\subsection{Simulations of phylogenetic sampling bias}

We tested whether the 77-taxon study sample could recover a period of increased
phenomic change. We simulated single-origin binary characters on the 765-taxon
phylogeny of \citet{blaimer2023} and reconstructed their rates before and after
pruning to the study taxa.

Within Apocrita, excluding its stem, the 0--80, 80--160 and 160--280~Ma
intervals contained 540, 177 and 17 eligible nodes. We randomly selected
5\%, 10\%, 30\%, 50\% or 80\% of nodes in each interval without replacement,
rounding to the nearest whole node (minimum one). Each selected node defined
one character: descendants were coded~1 and all other taxa~0. We generated
100 trials per combination, giving 1,500 paired datasets.

On both trees, we fitted each variable character with a homogeneous equal-rates
Mk model in \texttt{ontophylo} (development version 0.3.0), using five
optimization starts and the \texttt{maddfitz} root treatment
\citep{fitzjohn2009}. We generated 100 \texttt{phytools} stochastic maps per
character per trial \citep{revell2024}, holding the fitted model fixed.
Characters made constant by pruning were omitted.

We summed mapped changes across characters on each branch and divided by its
duration. We then averaged these combined branch rates across lineages present
at each age to obtain marginal curves at 1~Myr intervals. We compared all
simulated characters and, separately, the same surviving characters on both
trees.

To summarize recovery of recent change, we compared average rates in 0--80 and
80--160~Ma, weighting branches by their duration within each interval. We
counted trials with a higher recent mean and those whose conditional 95\%
mapping interval for the difference lay above zero. These counts describe
recovery under the imposed changes; they are not calibrated detection power.

\paragraph{Recovery of recent change.}
At the lowest intensity (5\% of recent nodes; 27 characters), only 4.94
characters remained variable on average: approximately 82\% were lost.
Nevertheless, the study tree recovered a higher recent rate in 92 of 100
trials, with a median 1.80-fold increase over 80--160~Ma. The conditional 95\%
mapping interval supported an increase in 86 trials. Recovery reached
98 of 100 trials at 10\% intensity and all 300 trials at 30--80\%.
Thus, a recent increase remained recoverable in almost all trials under this
design (490/500; 98\%), including most trials with sparse change.
